% DOT_IMPORT 08 8.14至8.15_校订续译工程
\egas{8.14}{推论（原稿陈述及待核处）}
结合8.13、8.5和8.6，按通常方式可得如下推论。\transnote{此为原作者的引导句；底本未另列完整推导，本译不把它当作已经进行独立验证的声明。}

假设 $f\colon X\to P$ 非分歧，且 $X$ 没有孤立点。〔字迹不清〕$n$。\footnote{原脚注49：“询问 A.G.”。}\textsuperscript{,}\transnote{假设末尾的残字 $n$ 及缺文照录，无法确定其数学含义，不据此补造附加假设。}\textsuperscript{,}\transnote{“直纹的”是原稿的暂用术语，不能径同于通常直纹簇概念；底本8.17专门指出二者不一致。本条沿用第十七版译稿，完整保存原陈述的充要方向。第十八版已用第108批回校稿替换8.13，并明示回校采用的新增假设；这项替换不等于本条8.14及其依赖关系已经重新审定，故这里仍不提供未经核实的补证。}
\begin{enumerate}
\item[(a)] 假设 $X$ 在 $k$ 上可分。\transnote{底本用词为 \emph{separable}。此处按原用语译为“可分”，不擅自改成“光滑”或“几何正规”；该旧用语在全书中的精确定义仍须结合前文核定。}
为使那些满足 $\phi_\xi$ 为 $X_{k(\xi)}$-正则、且 $Y_\xi$ 可分的点 $\xi\in P^\vee$ 所成集合，其补集的余维至少为 $2$，充要条件是：$X$ 的每个一维不可约分支 $X_i$ 都由 $X$ 的光滑点组成，并且相对于 $f$ 是“直纹的”；此外，对于 $X$ 的每个满足 $\dim_x X\ge2$ 的闭点 $x$，都有
\transnote{底本将 $\phi_\xi$ 误排成 $\phi_x i$，按本节一贯记号订正。深度条件印作 $\operatorname{prof}_x X\ge2\operatorname{prof}_x X$，重复了式末的 $\operatorname{prof}_x X$；正文删除该重复。依据是本句随后关于几何正规概形的说明，以及(b)的平行深度条件。按字面把它解释成乘积会使正深度点不满足条件，与原文紧随的说明不相容。此为排印校订，不是宣称已验证整个充要条件。$\operatorname{prof}$ 表示深度。}
\[
\operatorname{prof}_x X\ge2.
\]

（若 $X$ 几何正规，且所有不可约分支的维数均至少为 $2$，这些条件便自动满足。）
\item[(b)] 假设 $X$ 几何正规。为使那些满足 $\phi_\xi$ 为 $X_{k(\xi)}$-正则、且 $Y_\xi$ 几何正规的点 $\xi\in P^\vee$ 所成集合，其补集的余维至少为 $2$，充要条件是：$X$ 的每个维数不超过 $2$ 的不可约分支 $X_i$ 都由 $X$ 的光滑点组成，并且相对于 $f$ 是“直纹的”；此外，对于 $X$ 的每个满足 $\dim_x X\ge3$ 的闭点 $x$，都有
\[
\operatorname{prof}_x X\ge3.
\]
\end{enumerate}

\egas{8.15}{备注}
在8.6、8.13和8.14中，我们对 $X$ 作出 $(S_k)$（相应地，$(R_k)$；又相应地，可分、几何正规）这一假设，而希望将同一性质作为超平面截面的结论；唯一可能的例外，是 $\xi$ 属于某个余维至少为 $2$ 的例外集合。

这并不限制一般性；说实话，最好去掉这一预先假设。因为借助第3.4段和5.12的结果，可以立即看出：若 $X$ 不满足所讨论的假设，那么（由第5段），如果存在一个使该假设失效的闭点 $x$，则对于每个满足以下条件的 $\xi$：
\begin{enumerate}
\item $\phi_\xi$ 为 $X_{k(\xi)}$-正则（这个条件只排除一个余维为〔字迹不清〕的集合）；
\item $x\in Y_\xi$（这个条件描述一个余维恰为 $1$ 的集合），
\end{enumerate}
$Y_\xi$ 在 $x$ 处都不满足上述假设。于是，例外集合 $Z\subset P^\vee$ 的余维是 $1$，而不是 $2$。

（对于 $E_k$ 的情形，我也许说得有些过头了；在那里仍需要某种条件：$(S_1)$，也许还需要等维性，也许……）\transnote{这一保留意见为原文的一部分，必须与前面“不限制一般性”的断言一并阅读。底本的 $E_k$ 记号及缺失的余维数字均未恢复；不能把本段改写成无条件成立的去假设定理。}
\subsection*{补记：重新考虑（Remorse）}
在8.13和8.14中，只须假设 $f\colon X\to P$ 在 $X$ 的光滑点处非分歧；对于充分性部分，更只须它在 $X$ 的某个开子集 $X'$ 上非分歧，且其补集的余维至少为 $k+1$。
\transnote{底本将这段另题为“Remorse”，这里译作作者重新考虑前述假设的补记。它没有展开验证，且 $k$ 在8.14不同性质中的取值未在此重述；本译保留原范围和充分性限定，不据此全面替换前面条目的假设。}

% DOT_IMPORT 09 8.16至8.17_校订续译工程
\egas{8.16}{命题（原稿陈述）}
假设 $f\colon X\to P$ 在一个开子集上非分歧，该开子集的补集余维至少为 $2$；假设 $X$ 几何正规，且在其闭点处的深度至少为 $3$；最后，假设 $X$ 几何整且在 $k$ 上固有。那么，使得 $Y_\xi$ 几何正规、几何整且维数等于 $\dim X-1$ 的点 $\xi\in P^\vee$ 所成集合（是可构造的，并且）其补集的余维至少为 $2$。

\paragraph{原稿证明。}
事实上，由8.14(b)，对于性质“$Y_\xi$ 几何正规且维数为 $\dim X-1$”，上述结论成立（这一维数性质表达的是 $\phi_\xi$ 为 $X_{k(\xi)}$-正则）。\transnote{8.14前后的假设调整见8.15后的作者补记；其原文缺失和依赖问题仍未全部恢复。此处保留作者引用，不把引用本身作为此处译校已独立验证的证据。}

再由6.1，所有这些 $Y_\xi$ 都几何连通。由于 $Y_\xi$ 几何正规，它几何整当且仅当它几何连通，由此得到证明。
\egas{8.17}{备注}
\begin{enumerate}
\item[(a)] 8.16的假设蕴涵 $\dim X\ge3$。或许只要 $X$ 几何不可约且 $\dim f(X)\ge3$（不要求正规性和非分歧性），使得 $Y_\xi$ 几何不可约的点 $\xi$ 所成集合的补集就具有至少为 $2$ 的余维。\transnote{底本在此夹有法文“dès que”的转写痕迹。正文译为“只要”，并保留作者“It is possible”的推测语气，不将其升级为定理。}无论如何，我们可以证明这个补集不包含任何超平面（见下文）。\transnote{原句代词“it”依上下文指例外补集；此处明确所指，不改变其断言范围。}
\item[(b)] 如果去掉闭点 $x$ 处 $\operatorname{prof}_xX\ge3$ 的假设，8.16的结论便不成立。例如，对 $P^3$ 中的一个非奇异二次曲面 $X$，结论就不成立。〔字迹不清〕\footnote{原脚注50：“询问 A.G.”。}这些超平面截面是可约的（事实上由两条相交直线组成），因而形成一个二维族，在 $P^\vee$ 中的余维为 $1$（事实上，该二次曲面的对偶是在对偶空间中由对偶二次型给出的二次曲面……）。\transnote{底本此处印作“在 $P$ 中”，但所计数的是超平面的参数族，且紧随文字明确提到对偶空间，故订为 $P^\vee$。缺文处未补入“相切”等字样；这里的“这些超平面截面”只承接原稿所讨论的可约截面，不是说该曲面的所有超平面截面都可约。}
\end{enumerate}

对于射影空间中的非奇异曲面，依下一节（或下一条）的讨论，这种情形却应当视为例外情形。假设〔字迹不清〕在 $k$ 上整且固有，并且 $f$ 是浸入。那么由6.1、8.8及8.14可得：如果 $Y^{\mathrm{sing}}\to P^\vee$ 不是泛有限且不可分的，〔原稿接“one of the (???)”，句子残缺〕那些使得 $Y_\xi$ 在 $k(\xi)$ 上可分、并且至多有两个不可约〔字迹不清〕的 $\xi\in P^\vee$，〔句子残缺〕一个余维至少为 $2$ 的补集。\transnote{本段从假设到结论均有缺字和断句。“not generically finite and inseparable”的否定范围亦不明确；正文保留其不确定性，不自行判定应否定合取，或仅否定“泛有限”。不能由这些残句恢复一个确定的充分条件或结论。8.8亦尚缺完整假设，故不补造推论。“separable”沿用前批译为“可分”。}

现在我们要更仔细地考察曲面的情形（从超平面截面的不可约性这一观点看，曲线的情形显然不成为问题）。\transnote{“显然”是原作者的表述，底本此处未给证明；此处译校不另加未经核对的论证。}

（注意：我惊恐地发现，按我一直使用“直纹的”一词时所取的意义，二次曲面竟没有资格称为直纹的。这与前人的用法不一致，有必要为这里使用的概念另造一个更恰当的词。）\transnote{这是作者对暂用术语的自我修正，不是关于通常意义下二次曲面是否直纹的否定。}

% DOT_IMPORT 10 8.18_校订续译工程
\egas{8.18}{命题（原稿陈述，含待核处）}
假设 $k$ 是代数闭域，$X$ 是维数 $d\ge2$ 的整概形（分别地，整且正规），并且在 $k$ 上固有。设 $T$ 是 $X$ 的有限闭子集，使得 $X-T$ 光滑，并设 $f|_{X-T}$ 非分歧。要使那些满足 $Y_\xi$ 几何不可约（分别地，几何整）且“维数为 $1$”的点 $\xi\in P^\vee$ 所成集合的补集余维至少为 $2$，原稿称下列条件为充要条件。\transnote{底本此处确实写“dimension 1”，而(a)、(b)和随后的证明均用 $d-1$。末尾附记又讨论是否应限于 $d=2$，故不能仅凭上下文断定是漏写 $d-$ 还是原拟限定曲面。正文明确引述这一未决读法，不把它作为已校定的普遍定理。这里设 $d=\dim X$，只是明示原稿随后使用的记号。}
\begin{enumerate}
\item[(a)] 对每个 $x\in T$，存在一个经过 $x$、维数为 $d-1$ 且不可约的超平面截面 $Y_\xi$，其中 $\xi\in P^\vee(k)$。\transnote{底本(a)印作 $P(k)$，但 $\xi$ 参数化超平面，命题首句及(b)均使用对偶空间；这里统一为 $P^\vee(k)$。下文作为参数空间出现的 $P$ 同样统一为 $P^\vee$。}
\item[(b)] $X'=X-T$ 对于 $f$ 是“直纹的”（原稿如此），或者存在 $X'$ 的一个超平面截面 $Y'_{\xi_0}$，其中 $\xi_0\in P^\vee(k)$，其维数为 $d-1$，具有〔“非”字待核〕奇异性，且不可约。\footnote{原脚注51：“字迹不清，询问 A.G.”。}\transnote{底本写“(non???) singular”。究竟要求奇异还是非奇异，现有文本不能确定；这里不择其一。“直纹的”是作者暂用术语，应连同8.17后关于该术语的自我修正一起理解。}
\end{enumerate}
\subsection*{原稿证明}
先假设 $X$ 几何正规。由8.14(a)，我们已经知道可以在 $P^\vee$ 中找到一个余维至少为 $2$ 的闭子集 $Z'$，使得对 $\xi\in P^\vee-Z'$，$Y_\xi$ 在 $k(\xi)$ 上可分且维数为 $d-1$。对于这样的 $\xi$，$Y_\xi$ 几何不可约与几何整是等价的，所以8.18所考察的两个问题〔底本夹注：“respé et non respé”，法文第50页，问号〕便是等价的。

另一方面，由5.6，使 $Y_\xi$ 几何整且维数为 $d-1$ 的点 $\xi\in P^\vee$ 所成的集合 $U$ 是开的。这里维数假设表达的是 $\phi_\xi$ 为 $X_{k(\xi)}$-正则。

我们将给出一个包含于 $U$ 的非空开子集 $P^\vee-Z$，其定义是直接的：取 $Z$ 为 $g(Y'^{\mathrm{sing}})$ 与那些由 $f(x)$、$x\in T$ 确定的 $P^\vee$ 中的超平面 $H_x$ 的并。\transnote{底本此处像集未标闭包；正文不增添横线。$X\in T$ 的大写 $X$ 订为点变量 $x\in T$，与同句 $f(x)$、$H_x$ 及(a)一致。}对于 $\xi\in P^\vee-Z$，$Y_\xi$ 光滑且维数为 $d-1$；又由6.1，它几何连通，因而几何整。

因此，我们要表明（解释、证明），$Z$ 的每一个余维为 $1$ 的不可约分支都与开集 $U$ 相交。然而，这些不可约分支就是各 $H_x$〔它们可能重复，但这并不要紧〕，以及当 $g(Y'^{\mathrm{sing}})$ 确有余维 $1$ 时的这个集合；后者对应于 $X'$ 对于 $f$“不是直纹的”这一情形。

（注意：我们使用了 $Y'^{\mathrm{sing}}$ 的不可约性。另一方面，要使后一个集合与开集 $U$ 相交，原稿称充要条件是 $g(\overline{Y'^{\mathrm{sing}}})$ 与 $U$ 相交，并在这里插入“它包含一个开稠密集”的说明。）\transnote{底本这第三处横线位于 $Y'^{\mathrm{sing}}$ 上方，与前两处记号不同；不能未经原稿核对将其改为像的闭包。底本这句括号没有正常闭合，“which contains an open and dense set”的指代不够清楚。正文保留其内容而不擅自补入一个新的拓扑断言。按通常可构造像的论证，可以考虑像在其闭包中包含开稠密子集，但这里不把这种解释认定为手稿缺文的恢复。}

这就证明了此情形下的8.18。如果不假设 $X$ 正规，则对 $X$ 的正规化应用前面的结果；论证是直接的，我在这里不写出细节。\transnote{这是作者的省略声明，不是此处译校完成了正规化步骤的全部假设核验；尤其原稿(b)的奇异性条件仍待核。}

注意：在〔respé，即“分别地”〕的情形，8.18包含在8.16中；更确切地说，要除去 $d=2$ 的情形。至于〔non-respé，即非“分别地”〕的情形，或许最好“不要求 $d=2$，而不只是 $d\ge2$”……\transnote{最后这句话按底本“not to require $d=2$ and not only $d\ge2$”保留，其否定与比较关系在当前转写中不清楚。不能据此确定应增加还是删除曲面限制，亦不能据此前文维数 $1$ 的歧义宣称已修复。}

% DOT_IMPORT 11 8.18后线性系讨论_第一批_工程
\subsection*{（8.18）之后的线性系讨论}
还须把8.18的条件(a)、(b)解释清楚（明确写出）。\transnote{底本印作“1.18”，但紧接8.18并明确指向该命题的(a)、(b)，此处订为8.18。}这使我们一般地考察如下情形。假定 $X$ 在 $k$ 上几何不可约，并给定一个超平面参数空间中的线性子簇 $L$。\transnote{底本此处写“a linear subariety $L$ of $P$”，随后又用 $L$ 表示被这些超平面包含的线性子簇，并在后文写作 $L^0$。为区分两个不同空间，正文将参数线性系称为 $L$，其基线性子空间称为 $L^0\subset P$；此举不补定后面关于 $F^0$ 的未明约定。}这对应于研究经过给定点 $x$，或在给定光滑点与 $X$ 相切的超平面截面。因此，$L$ 由那些包含 $P$ 中某个线性子簇 $L^0$ 的超平面组成；在上述两种情形中，$L^0$ 分别是一个点，或 $X$ 在光滑点的切空间的像。我们提出如下问题〔底本夹注：“de ravoi”，读法待核〕：对于 $L$ 的泛点（因而对于 $L$ 的某个非空开子集的所有点），$Y$ 是否几何不可约，且
\[
\dim Y=\dim X-1?
\]
这是 Bertini 定理的一个变式，应当放在第3节中〔底本夹注：“j(devrait figurer) 51”〕，并且完全用同一方法处理〔或者，愿意的话，可以将它归结到那里；底本法文夹注“ou, si on veut, s'y ramène”，法文稿第51页〕。
维数问题可以简单地表述为 $f(X)\not\subset L^0$，也就是说，
\[
X'=f^{-1}(P-L^0)
\]
是 $X$ 的稠密开子集。\transnote{底本写成“$f(X)\notin L-0$”，下一句立即给出上述逆像稠密条件；这里订正为像不包含于中心 $L^0$，不是把像当成一个点。}
令 $Q$ 为原稿所谓“经过 $L^0$ 的超平面的射影空间”〔注意：如果 $L^0$ 由 $E$ 的向量子空间 $F^0$ 定义，原稿写 $Q=P(F^0)$〕，并考虑典范态射（由底本所记 $F^-\to E$ 导出，参见第II章）：\transnote{这一句的射影对偶约定尚不清楚：原稿同时把 $Q$ 称为超平面参数空间，又在下面使用 $L\simeq Q^\vee$，二者不能不加解释地混同。$F^-$ 的上标也与前面的 $F^0$ 不同。此处译校保留这些写法并标待核，不据此猜补向量空间的商、对偶或映射。}
\[
u:P-L^0\longrightarrow Q.
\]
再考虑\transnote{底本第二条显示公式写 $f^{-1}(P-L)$，但前面的 $X'$ 与 $u$ 的定义域均为中心 $L^0$ 的补集；此处统一为 $L^0$。$f'$ 指 $f$ 在 $X'$ 上的限制。}
\[
g=uf':f^{-1}(P-L^0)=X'\longrightarrow Q.
\]

于是 $L\simeq Q^\vee$，而 $X'_\xi$（$\xi\in L$）这一族，正是相对于态射 $g$ 的超平面截面族。
另一方面，原稿说可以立即看出：对于每个“足够一般的”$\xi\in L$，“$X'$ 在 $X$ 中稠密”，从而“$X'$ 几何不可约，当且仅当 $X$ 几何不可约”。\transnote{底本在这句话中未写 $\xi$ 下标，但上句正在讨论截面族 $X'_\xi$。可能本意是相应截面中的稠密性，而不是重复全空间 $X'\subset X$；因原稿未明写，正文引述其现存记号，不擅自补成对所有截面的断言。}
在此基础上，Bertini--Zariski 定理说明，只要
\[
\dim g(X')\ge2,
\]
便得到所需的不可约性结论。

（严格说来，原稿提出可将3.1的逆命题写成如下形式：如果 $X$ 几何不可约，则 $Y$ 几何不可约，当且仅当以下两种情况之一成立：或者 $\dim f(X)=2$；或者 $\dim f(X)=1$，且 $f(X)$ 包含在一条定义于 $k$ 上的直线 $D$ 中，并且 $X\to D$ 的泛纤维几何不可约。）\transnote{底本确实为“$\dim f(X)=2$”，不是“$\ge2$”；它与紧前一句所用的充分条件存在明显张力。这里引述作者拟议的逆命题并标待核，不能把“当且仅当”当成已核实的定理，也不在缺少完整假设和证明时擅自补成一个修订逆命题。}
原稿接着说，这也使得我们在当前带有 $L$ 的版本中，得到 $Y_\xi$ 几何不可约的充要条件，其中 $\xi$ 是 $L$ 的泛点。\transnote{这句话依赖上面尚待核的逆命题；忠实译出不表示已验证该充要条件。}

% DOT_IMPORT 12 第8节末域论讨论_工程
\subsection*{域论与维数的补充讨论}
从〔底本夹注：“(cunutesque?) Fr”，读法待核〕的观点出发，用域论的语言，可以如下用超越次数来表述这一条件。选择一个既不包含 $L^0$、也不包含 $X$ 的“无穷远超平面”，然后在其补集中讨论，也就是说，实质上在一个仿射型概形上讨论。\transnote{这里“包含 $X$”沿用底本措辞；在给定态射 $f:X\to P$ 的设置下，应按其像 $f(X)$ 理解，而不是假定 $f$ 已是嵌入。}
取在 $L^0$ 上消失的线性形式空间的一组基，记为 $T_1,\ldots,T_p$，其中 $p=\operatorname{codim}(L^0,P)$；考虑它们在 $X$ 的分式域 $K$ 中的拉回 $t_1,\ldots,t_p$（这里假定 $X$ 整）。\transnote{底本在这句写“vanishing on $L$”，但本段中心为 $L^0$，且紧接着指定 $p=\operatorname{codim}(L^0,P)$；这里统一中心记号为 $L^0$，区别于前段的超平面参数线性系 $L$。在所选仿射开集上，线性形式经平凡化成为函数，再拉回到函数域。}
至少有一个 $t_i$ 不为零，不妨设为 $t_1$。于是令
\[
a_1=t_2/t_1,\quad\ldots,\quad a_{p-1}=t_p/t_1.
\]
那么 $\dim g(X')$ 正是子域 $k(a_1,\ldots,a_{p-1})\subset K$ 对 $k$ 的超越次数。\transnote{底本写为“$K(a_1,\ldots,a_{p-1})\subset K$”。由于 $a_i\in K$，该域就是 $K$，一般不能给出像的维数。正文订为 $k(a_1,\ldots,a_{p-1})$：在投影的首坐标非零的仿射图中，像的坐标环由这些坐标比生成，其分式域即此子域，故像闭包的维数等于它对 $k$ 的超越次数。这一订正仅针对本句，并不恢复前批关于 $Q$ 的全部对偶约定。}
因此，若该超越次数至少为2，便符合要求。
若它等于1，则原稿要求：在 $k$ 上，$f(X)$ 包含于 $P$ 的某个包含 $L^0$、且维数至多为1的线性子簇中，并且 $g:X'\to g(X')$ 的泛纤维应当几何不可约。\transnote{“dimension at most one”是底本跨第35、36页的现存文字，不能直接当成已核实判据。若 $\dim L^0>1$，这样的线性子簇根本不能包含 $L^0$；而从中心 $L^0$ 投影后像落在直线的条件，通常对应于其逆像落在更高维的线性子空间。这里只明确引述原稿要求，未猜补可能漏失的相对维数或改写为未经证实的充要条件。}

设 $L^0$ 的维数为 $q$，则 $u:P-L^0\to Q$ 的纤维维数为 $q+1$。\transnote{底本先写“dimension $g$”，下一行及后续却均用 $q$；此处将维数变量统一为 $q$，不与态射 $g$ 混同。这里采用前文由中心 $L^0$ 给出的线性投影。}
将该投影限制到 $f(X)\setminus L^0$，其纤维维数至多为 $q+1$，因而\transnote{底本写“so that those of $g$ are $\dim\le q+1$”。一般 $g=u\circ f'$ 的纤维还包含 $f$ 自身的纤维，不能仅由 $u$ 的纤维界直接得到这一说法。正文改用投影在像 $f(X)\setminus L^0$ 上的限制：其纤维是 $u$ 的纤维中的子集，维数界及随后关于像维数的不等式成立。若另知 $f$ 拟有限，原稿对 $g$ 的纤维界亦可成立；此处不擅自新增该假设。像的维数按其闭包的维数理解，并沿用 $X'=f^{-1}(P-L^0)$ 稠密的前提。}
\[
\dim g(X')\ge \dim f(X)-(q+1).
\]

所以，若 $\dim f(X)\ge q+3$，便满足关于 $g(X')$ 的维数条件。当 $q=0$ 时，就得到8.17(a)中所指出的事实。
回到8.18的条件，原稿指出：对于一个 $x\in T$，只要 $x$ 不是“相对于 $f$ 在 $x$ 处为锥形的”，条件(a)便满足，这里的意义是显然的。\transnote{底本为“provided $x$ is not ‘conical at $x$ relative to $f$’”。首个小写 $x$ 可能本应为 $X$，也可能是在描述点 $x$ 为相对于 $f$ 的锥点；对象和定义均未明，正文保留现存措辞，不确定恢复。原稿未在本段给出精确定义或完整论证，故这句不作为已经独立验证的判据。}
也许最好把这些最后的 Bertini 式讨论放到下一节中……更换射影浸入。

% DOT_IMPORT 13 9.1至9.2_校订续译_工程
\egasection{§9\quad 更换射影嵌入}\source{\texttt{EGAVaiello.pdf} 第36--38页}
\egap{9.1}
对每个整数 $n>0$，令
\[
P(n)=\mathbb P(\operatorname{Sym}^{n}(E)).
\]
有一个显然的浸入 $u_n:P\to P(n)$，因为在 $S$ 的每个仿射开集上，$\mathcal O_P(n)$ 都由其截面生成，并且
\[
p_*\mathcal O_P(n)\simeq\operatorname{Sym}^{n}(E),
\]
其中 $p:P\to S$ 是投影。\footnote{原脚注52：“无法辨认。”底本此处印为 $p_*(0_p(n))\operatorname{Sym}^{n}(E)$，缺关系符，前式还将 $E$ 印作小写 $e$。译校说明：正文统一记号，并依射影丛的商约定写出标准典范同构；这是数学上可核实的整理，不声称恢复了手稿原字迹。由评估满射 $p^*\operatorname{Sym}^{n}(E)\twoheadrightarrow\mathcal O_P(n)$ 得到相对 Veronese 浸入；在 $E$ 平凡化的开集上，其坐标为全部 $n$ 次单项式。}
若 $f:X\to P$ 是非分歧态射（分别地，是浸入），则 $u_n f:X\to P(n)$ 也如此。有时，研究 $X$ 时用 $u_n f$ 代替 $f$ 是有利的，可以避免 $f$ 在某些方面所具有的十分特殊、且有时令人棘手的行为。
（这种特殊情形的一个例子是8.12(b)中所指出的〔底本原有“sic”标记〕情形：$Y^{\mathrm{sing}}\to P^\vee$ 的像具有维数 $r-1$，但产生的域扩张不可分。另一个例子由 $\mathbb P^3$ 中的二次曲面给出：其所有奇异超平面截面都几何可约，尽管 $X$ 本身几何不可约。）\transnote{这里按原文保留二次曲面这一例子的概括措辞，未扩大为任意特征、任意退化二次曲面的无条件断言。前文几何普通奇点的术语约定沿用原稿；不将“几何可约”与“非约化”混同。}

\medskip
\egas{9.2}{命题（原稿陈述，所引条件待核）}
设 $S=\operatorname{Spec}(k)$，$X$ 在 $k$ 上光滑，且 $f:X\to P$ 非分歧。取 $n\ge2$，并考虑 $f_n=u_n f$。原稿断言：$f_n:X\to P(n)$ 满足8.8中的等价条件；特别地，在超平面参数空间 $P(n)^\vee$ 中，除去一个余维至少为2的集合后，对于其中的点 $\xi$，相应的超平面截面 $Y_\xi$ 或者光滑，或者只有有限个非光滑点，且这些点都是几何普通奇点。若 $f$ 是浸入，则这样的奇点至多只有一个，并且它在 $k(\xi)$ 上有理。\transnote{底本写 $\xi\in P(n)$。由于 $\xi$ 参数化 $P(n)$ 中的超平面，正文订为 $P(n)^\vee$，与9.1所用 $P^\vee$ 及全篇的超平面参数约定一致。此为参数空间记号订正，并非补全8.8条件。底本此处未给出完整证明；下文仅有关于8.8(ii bis)的核验提示。}

\medskip
\noindent\textbf{附记（N.B.）.}
应当以一种不排除 $f$ 并非浸入这一情形的方式来陈述8.8。在8.8的条件(ii bis)下，核验本质上是平凡的。\transnote{“核验本质上是平凡的”为原稿评语，不是本次独立审校的结论。现有底本缺失8.8的标题和完整假设，虽残存含(ii bis)的条件列表，仍无法仅凭这句完成核验，也不能据此声称已恢复缺文。}
无疑，应当在9.1中明确说明：$X$ 相对于 $f_n$ 的超平面截面，无非是用次数为 $n$ 的超曲面代替超平面所得的 $X$ 的“截面”。

% DOT_IMPORT 14 9.3_校订续译_工程

\egas{9.3}{命题（原稿陈述；(c)项措辞待核）}
设 $X$ 几何不可约，且 $\dim f(X)\ge2$；取 $x\in X(k)$ 及 $n\ge2$。
\begin{enumerate}
\item[(a)] 考虑 $X$ 相对于 $f_n=u_n f$ 的、经过 $x$ 的超平面截面的线性族。其泛成员所定义的 $Y^{(n)}_\xi$ 几何不可约。
\item[(b)] 设 $L$ 是 $P$ 中经过 $f(x)$ 而不包含 $f(X)$ 的线性子簇。考虑 $X$ 相对于 $f_n$ 的如下超平面截面线性族：它由 $P(n)$ 中“在 $x$ 处与 $L$ 相切”的超平面定义，即由 $P$ 上的 $n$ 次齐次式定义，这些式子限制到 $L$ 后在 $x$ 处至少二阶为零。其泛成员 $Y^{(n)}_\xi$ 几何不可约。\transnote{这里在 $L$ 上的点应理解为 $f(x)$；沿用底本“在 $x$ 处”的简写。条件是限制到 $L$ 后在该点至少二阶消失，不是沿整个 $L$ 的理想平方条件。}
\item[(c)] 设 $X$ 在 $x$ 处光滑，且 $n\ge3$〔底本接着写：“where $f(X)$ is not contained in a plane defined over $k$”，即“其中 $f(X)$ 不包含于一个在 $k$ 上定义的平面中”；连接关系待核〕。考虑 $X$ 相对于 $f_n$ 的、“在 $x$ 处与 $X$ 相切”的超平面截面 $Y^{(n)}_\xi$ 族。其泛成员所定义的 $Y^{(n)}_\xi$ 几何不可约。\transnote{底本确实印为“and that $n\ge3$ where…”，不是可确认的“or”。证明随后区分切空间像的维数，并在二维情形使用三次式、较高维情形使用二次式，提示条件可能另有原意；此处译校不据此猜改连接词，不声称恢复手稿。}
\end{enumerate}

\noindent\textbf{证明（原稿）.}
利用上一节末尾的判据，证明本质上是平凡的。\transnote{此句为原稿评语，不是本次独立验证的结论。后文照译论证，并单列可核实的记号订正与未决问题。}
在 $P$ 中取包含 $f(x)$ 的仿射坐标开集，并在其逆像上考虑 $X$，\transnote{底本写“取一个包含 $f(X)$ 的 $P$ 的仿射模型”，一般不能要求整幅射影像包含于一个仿射开集。这里改为经过 $f(x)$ 的仿射坐标邻域，其逆像非空，因 $X$ 不可约而稠密，足以定义下文函数域中的坐标。此局部化并不自行完成对全截面的不可约性论证，后者仍使用上一节判据。}
在(a)中，问题归结为：在坐标 $T_1,T_2,\ldots,T_r$ 中找到三个次数不超过2的多项式 $P,Q,R$，使得
\[
\frac{Q(t)}{P(t)},\qquad \frac{R(t)}{P(t)}
\]
在 $K$ 中于 $k$ 上代数无关。这里 $K$ 是 $X$ 的函数域，$t=(t_1,\ldots,t_r)$ 是由 $T_i$ 定义的 $K$ 中元素系。\transnote{如 $X$ 未约化，函数域按其不可约约化概形理解。本段坐标可取以 $f(x)$ 为原点，从而后续坐标式在该点消失；这是下文所用的坐标约定。多项式 $P$ 与射影空间 $P$ 同名，沿用底本。}
在(b)中，还要求 $P,Q,R$ 限制到 $L$ 后在 $x$ 处至少二阶消失；此外可设 $L$ 的方程为\transnote{底本这里写“vanish to order two at least on $L$”。若解作 $P,Q,R\in I_L^2$，便与下文 $P=T_1$ 冲突。按(b)项原有括号中的定义，正文明确为限制到 $L$ 后在点上二阶消失；下文 $P,Q,R$ 限制到 $L$ 恒为零，确实满足这一条件。}
\[
T_1=\cdots=T_s=0.
\]

最后，(c)中的情形相同，只是 $L$ 为 $X$ 在 $x$ 处切空间的像，并且必要时允许取次数为3的 $P,Q,R$，也就是说，可以稍稍放宽次数。

假设 $\dim f(X)\ge2$ 意味着 $k(t_1,\ldots,t_r)$ 在 $k$ 上的超越次数至少为2。因此，可以找到两个代数无关的坐标函数，不妨记为 $t_1,t_2$。\transnote{底本印作 $K(t_1,\ldots,t_r)$，但所有 $t_i$ 已在 $K$ 中，故该式等于 $K$，不能表达像的维数；正文订为 $k(t_1,\ldots,t_r)$。}
于是在(a)中取
\[
P=T_1,\qquad Q=T_1^2,\qquad R=T_1T_2.
\]
在(b)中也类似地这样取。注意，此时可选择使 $T_1$ 在 $L$ 上为零的 $t_1$，这是因为 $f(X)\not\subset L$。这意味着存在 $1\le i\le s$ 使 $t_i\ne0$；由于 $t_i$ 在 $x$ 处为零，它不可能是常数，因而不在 $k$ 上代数。\footnote{原脚注53：“〔英译者〕由译者补入。”该脚注在底本中附于“不在 $k$ 上代数”之后。}\transnote{底本印为 $f(x)\not\subset L^*$，若星号不表示另一个对象，则与(b)中的 $f(x)\in L$ 相抵触；正文按(b)假设改为 $f(X)\not\subset L$，不为无说明的星号添加对象。底本先说存在下标 $i$，随后却写 $t_s$，正文统一为所选 $t_i$。原文紧接着允许设 $k$ 代数闭；在此假设下，函数域中代数于 $k$ 的元素属于 $k$，故所述推理成立。}
（附记：可以假设 $k$ 为代数闭域。）

除去 $f(X)$ 包含在 $X$ 于 $x$ 的切空间经 $f$ 所得的像 $L$ 中这一情形外，(c)由(b)推出。〔若 $\dim X=2$，这一情形确实是例外：一个二次曲面与其相切平面的交迹，一般由两条相交直线组成，因此不可约性不成立。〕\transnote{方括号为原稿所有。这里保留其例示性质，不扩大为任意特征或任意退化二次曲面的无条件结论。}
但为处理这一情形，在上面明确给出的 $P,Q,R$ 中，可以将 $X$ 换成 $L$ 本身，而在 $L$ 上问题的解是平凡的。\transnote{这是原稿的归约断言。若仅知 $f(X)\subset L$，在 $L$ 上代数无关的函数限制到 $f(X)$ 后不必仍代数无关。因而此处译校不把这句视为已补全的论证；需结合原假设与切空间像的维数进一步核查，尤其不能单靠下面的多项式计算确认(c)项的全部范围。}
（若 $\dim L=2$，取
\[
P=T_r^2,\qquad Q=T_r^3,\qquad R=T_r^2T_{r-1},
\]
注意 $Q/P=T_r$ 与 $R/P=T_{r-1}$ 是在 $L$ 上线性无关的线性式，因而代数无关。若 $\dim L\ge3$，则 $T_{r-2},T_{r-1},T_r$ 在 $L$ 上线性无关，取
\[
P=T_r^2,\qquad Q=T_{r-1}^2,\qquad R=T_{r-2}^2.
\]
）\transnote{底本在这段后缺右括号，正文仅补排版闭合。所取坐标在 $L$ 上的线性无关性按原文保留；在这样的线性空间坐标约定下，两比值的代数无关性可核验，但不自动解决上一条译注所述限制到像的问题。}

% DOT_IMPORT 15 9.4至10.1_校订续译_工程
\subsection*{（9.4）的编排提示与（9.5）的残文}
结合8.18，我们得到一个推论9.4。（〔英译者〕待陈述。）

最后，还须将后者与9.2结合，以便在“优良情形”下得到一个汇总性的定理。

\egas{9.5}{定理残文（第52页或第55页字迹不清）}\footnote{原脚注54：“请询问A.G.，我未能理解。〔英译者〕”}

〔残文〕（若 $X$ 在 $k$ 上光滑、固有且几何整，且 $f:X\to P$ 非分歧，$X$ 的维数 $\ge2$（请询问Grothendieck），通过考虑结果 ???，当 $x\to p$ 是一个浸入时（$Y_\xi$ 的奇异点簇。\transnote{底本这段语法与括号均未闭合，实际印为小写 $x\to p$；此处保留而不据猜测补写定理。所谓第52或55页是底本自身的编排标记，不是此处译校PDF页码。}
\egasection{§10\quad 超平面截面束与簇的爆破的纤维化}\source{\texttt{EGAVaiello.pdf} 第38--41页}
\egap{10.1}
设 $Z$ 是 $P^\vee$ 中相对于一个可构造性质 $\mathcal P$ 的例外集，并假设 $Z$ 为可构造子集，$S=\Spec(k)$。例外集的余维按其闭包的余维理解。以下关于直线的判据以 $\dim P^\vee\geq1$ 为前提。\transnote{原稿直接讨论直线，未单列零维参数空间；零维射影空间不含直线，不能使用这里的直线存在性判据。}我们将看到（参见第12节，那里将补述一些无疑本应在前几节出现的内容），条件
\[
\codim(\overline Z,P^\vee)\geq2
\]
等价于：几何意义下充分一般的直线 $L$ 不与 $\overline Z$ 相交。对于闭集 $\overline Z$，只要存在一条与之不相交的几何直线便足够。\cmt{D17-E03}{底本在“$Z$”“或者甚至 $Z$”以及“应为 $Z$”处没有可辨的上划线，并称存在单条不交 $Z$ 的直线便足够；这对一般可构造集不成立。取射影平面中一条直线删去一点，经过所删点的另一条直线可避开该集合，而其闭包的余维仍为1。正文明确使用闭包；旧句及编辑批注完整保存于报告。底本脚注55为“请询问A.G.”。此订正不声称恢复了未见的法文手稿。}
若 $k$ 无限，上述余维条件又等价于存在一条定义于 $k$、不与 $\overline Z$ 相交的直线 $L\subset P^\vee$。\transnote{底本此句印为小写 $\bar z$，正文按同一例外集写作 $\overline Z$。闭集余维至少2时，避开它的直线构成 Grassmann 簇的非空开集；其仿射坐标图在无限域上有稠密的 $k$-有理点。反之，任何余维1的射影闭子集都与每条几何直线相交。这里区分几何直线与定义于 $k$ 的直线。}
我们把 $L$-预概形 $Y_L$ 称为由 $P^\vee$ 中直线 $L$ 定义的 $X$ 的一个超平面截面线性束（此定义对任意 $S$ 都有效）。

于是，将上述考虑与第8、9节的结果结合，首先在 $S$ 为一个无限基域的谱的情形下，便得到这样的束的存在性判据：所有纤维 $Y_\xi$（$\xi\in L$）都满足性质 $\mathcal P$。考虑到8.2，若对 $X$ 的每一个伴随素循环 $T$ 都有
\[
\dim f(T)>0,
\]
则可以要求束 $Y_L$ 在 $L$ 上平坦，办法是取性质
\[
\mathcal P'=\mathcal P+\text{$\phi_\xi$ 的正则性条件}.
\]
在 $S$ 任意的情形下，若 $k(s)$ 无限，而且已知 $Z$ 为闭集，则仍可按照7.1的方法，在 $S$ 的给定点 $s$ 的某个开邻域上构造这样的束。（第5节的结果以及 $X\to S$ 固有的假设，在各种不同情形下保证上述闭性。）

为了妥善处理此事，宜在作出这种一般性说明后，给出若干汇总性陈述，将前面的结果实际应用于若干这类性质（也包括模的性质）。至少，必须在这里用线性束的语言重新表述9.5；这是一项在几何应用中经常使用的事实。\transnote{此句是原稿的写作计划及评语，并不意味着底本中缺失的9.5已经补齐。}

% DOT_IMPORT 16 10.2_校订续译_工程
\egap{10.2}
借助极性对应，$P^\vee$ 中的一条直线 $L$ 对应于 $P$ 中余维为2的线性子簇 $L^0$（$S$ 任意）。置\transnote{底本此处印作“$T=x_P L^0$”，与随后零点概形的描述相结合，可明确读为沿 $X\to P$ 的逆像，即上述纤维积；这里补全纤维积符号并统一 $X$ 的大小写。}
\[
T=X\times_P L^0.
\]

还可以如下描述 $T$：$L$ 由 $E^\vee$ 的一个秩为2的局部自由商模定义；等价地，由 $E$ 的一个处处秩为2、局部为直和因子的子模 $F$ 定义。考虑复合同态
\[
F_X\longrightarrow E_X\longrightarrow\mathcal O_X(1),
\]
则 $T$ 正是这个复合同态的零点概形；换言之，它由理想 $J$ 定义，其中 $J$ 是相应同态（通过张量乘以 $\mathcal O_X(-1)$ 得到）
\[
F_X(-1)\longrightarrow\mathcal O_X
\]
的像。

假设这个同态是正则的，意思是：在局部写出 $F_X(-1)$ 的一个全序基，其在 $\mathcal O_X$ 中的像构成一个 $\mathcal O_X$-正则序列。这个条件与所选的基无关；也可以内在地表述为：
\[
F_X(-1)\otimes_{\mathcal O_X}(\mathcal O_X/J)
\longrightarrow J/J^2
\]
是同构，并且 $V(J)=T\to X$ 是正则浸入。\transnote{这里忠实保留原稿的正则性表述；在任意预概形上，所用“正则序列”的局部约定、是否仅在零点概形附近要求该条件，以及更换所选局部基后的不变性与上述内在表述的确切等价范围，应与作者有关正则浸入的一节统一核查。此处译校未独立证明这一等价。}
〔原稿附记：应当在某处阐明一般情形，即预概形 $X$ 上有一个同态 $G\to\mathcal O_X$，其中 $G$ 局部自由；例如可放在关于正则浸入的一节。〕于是有：

\par\noindent\textbf{定理（10.2）。}
在上述假设下，带有典范投影 $Y_L\to X$ 的线性束 $Y_L$，以唯一的方式在 $X$ 上同构于预概形 $X$ 沿中心 $T$ 的爆破。

为了理解这个定理的意义，宜在本节或本条开头指出：若 $S=\operatorname{Spec}(k)$，则对 $P^\vee$ 中一条“充分一般”的直线 $L$，上述正则性条件成立（参见第12节所指的补述，即关于5.3的补述）。\transnote{这是原稿对后文的预告，不应理解为无须约定定义域便可在任意有限域上选取一条所需的 $k$-有理直线；其精确范围仍须结合5.3与第12节核对。}
在下文构造本条开头所指出或预期的“优良”线性束时，我们可以要求所述束满足这一条件（它与下述条件属于同一类型，但并不相同：对每个 $\xi\in L$，要求 $\phi_\xi$ 是 $X_{k(\xi)}$-正则的）。应把所讨论的条件纳入拟作的汇总性陈述。

另一方面，实际上10.2只用于9.5的情形，因此，最好在10.2之后才用束的语言重新表述9.5，以便在该陈述中同时纳入束与一个爆破的同构（即给出便于适当引用的情形描述）。这样，对于无限域 $k$ 上每一个维数 $\ge2$ 的光滑、几何连通的射影概形 $X$，我们便得到一种方法，找出一个非空的光滑闭子概形，它在其每一点的余维均为2，使得爆破后的概形具有到 $\mathbf P^1$ 的一个纤维化：所有纤维均几何整，并且除至多有限个纤维外，其余纤维都光滑；那些例外纤维至多有一个几何奇异点，而且这样的点在相应参数点 $\xi$ 的剩余域 $k(\xi)$ 上有理，并且在几何意义下是普通奇点。\transnote{底本确实写“rational over $k$”。纤维在参数点 $\xi$ 的剩余域上定义，故相应唯一几何奇异点的有理性应表述为在 $k(\xi)$ 上有理，不能无条件说成在 $k$ 上有理；正文据此订正其基域。这整段仍是原稿拟议的应用，所引用的（9.5）在现有底本中残缺，译校不声称已补齐它或独立证明全部存在性。}

这说明，深入研究这种带奇异纤维的纤维化是重要的（这种研究在目前才刚刚开始）：它可在一定程度上把维数为 $d$ 的光滑射影簇的研究，归约为对维数为 $d-1$、容许有普通奇点的射影簇的研究（这些簇组成依赖一个参数的族）。\transnote{底本在此句有未闭合括号及语法不顺的“or”；这里依可辨的语义整理括号与语序，没有加入新的数学结论。}

陈述10.2或多或少是下面命题的直接推论；后者完全独立于超平面截面这一主题，无疑更适合放在另设的一段“正则浸入”中。\footnote{原脚注56：“请询问A.G.”}

〔底本编排批注〕命题10.3被划去了。请询问A.G.，这是否是他的本意。

% DOT_IMPORT 17 10.3_校订续译_工程
\egas{10.3}{命题（底本标记已划去；以下保存其文字）}
设 $X$ 为预概形，$G$ 为 $X$ 上的拟凝聚模，
\[
u:G\longrightarrow\mathcal O_X
\]
为一个同态，并置 $J=u(G)$、$T=V(J)$。设 $\widetilde X$ 是由 $X$ 沿 $T$ 爆破而得的预概形。\transnote{底本把爆破前后的预概形均排作 $X$。为避免把后面的闭浸入误读为原预概形 $X$ 的闭浸入，本译文统一以 $\widetilde X$ 表示该爆破；这仅是区分记号，不改变所指定的对象。}
另一方面，考虑 $P=\mathbf P_X(G)$、典范同态
\[
G_P\longrightarrow\mathcal O_P(1)
\]
及其核 $H$，于是有正合列
\[
0\longrightarrow H\longrightarrow G_P
\longrightarrow\mathcal O_P(1)\longrightarrow0.
\]
再考虑同态 $u_P:G_P\to\mathcal O_P$，以及拟凝聚理想\transnote{底本写作小写 $p=p(G)$，本译文将射影预概形的记号统一为 $P=\mathbf P_X(G)$，并相应统一下标。$G_P$ 表示 $G$ 在 $P$ 上的拉回；$u_P$ 为相应拉回同态。底本的“quadi-coherent”按“quasi-coherent”译出。}
\[
K=u_P(H)\subseteq\mathcal O_P.
\]

则 $\widetilde X$ 典范地同构于 $V(K)$ 的一个闭子概形。

若 $G$ 局部自由，且 $u$ 是“正则的”，则上述同构就是 $\widetilde X$ 与 $V(K)$ 本身的同构。在这种情形下，$H$ 是 $P$ 上的局部自由模，而同态
\[
H\longrightarrow\mathcal O_P
\]
（其〔底本此处缺文：预概形之“???”〕是 $\widetilde X$）也是正则的。\transnote{底本末句实际排作“$H\to\vartheta_p$ 9whose prescheme of ??? is $X$) is also regular”。由前文 $u_P$ 的限制可识别这里的同态 $H\to\mathcal O_P$；但括号内“???”缺失的具体词语未恢复，因此没有将猜测的“零点概形”补作原文。本条“局部自由”及“正则”的秩与局部约定仍需与10.2及作者有关正则浸入的论述核对，此处译校未独立验证命题的全部一般性。}

\subsection*{原稿紧随其后的论证说明}
第一个断言几乎是平凡的。第二个断言是一个不会造成困难的练习（我没有详细做过，心想你和我一样能够把它推导出来）。\transnote{“几乎是平凡的”及括号内的第一人称均为原稿文字，并非本译者对证明充分性的判断。现有底本此处没有详细证明，不能将这一说明当作已完成的独立审证。}

% DOT_IMPORT 18 10.4_校订续译_工程
\subsection*{（10.4）之前的说明}
若在10.2中取 $S=\operatorname{Spec}k$，并限于下述光滑曲线的情形，则所作的正则性假设等价于 $T=\varnothing$，因而 $Y_L\to X$ 是一个同构。\transnote{底本实际写作“$X$ is of dimension $\geq1$”。这一维数条件不能支持原结论：例如二维光滑仿射空间中由两个坐标给出的正则截面具有非空零点中心。对于下述光滑曲线，局部环是诺特环且维数至多1；非空中心处若存在长度2的正则列，则深度至少2，与深度不超过维数矛盾。反向采用10.2中允许空中心的正则截面（正则浸入）约定。故正文只写已核实的曲线情形，不把猜测的“$\leq1$”冒充已恢复的原稿。}因此，将此与9.2结合，得到：
\egas{10.4}{推论（（10.2）的推论）}
设 $X$ 是无限域 $k$ 上射影空间 $P$ 中的一条光滑、几何连通的曲线，且 $n\geq2$。则存在 $P$ 上的一个 $n$ 次形式的线性束，它定义一个态射
\[
X\longrightarrow\mathbf P^1,
\]
并具有下述性质：这个态射泛平展，次数为 $d$；对于 $\mathbf P^1$ 的每一个几何点 $s$，$X_s$ 或者在代数闭域 $k'=k(s)$ 上平展，或者在 $k'$ 上同构于 $d-2$ 个概形 $\operatorname{Spec}k'$ 与概形
\[
I_{k'}=\operatorname{Spec}\bigl(k'[t]/(t^2)\bigr)
\]
的不交和。\transnote{底本在“$d-2$ 个概形”处印作 $\operatorname{Spec}k$；这里描述的是 $k'$ 上的几何纤维，故按其底域更正为 $\operatorname{Spec}k'$。$k'=k(s)$ 沿用底本的几何点记号；这里的 $s$ 是取值于代数闭域的几何点，不能误作底空间任意点的未代数闭剩余域。}

用前辈们的语言来说：每个这样的几何纤维中至多有一个分歧点，而且它是“二次的”。\transnote{底本末句为“there is at most one point of ramification and it is ‘quadratic.’”这里“至多一个”承接前句，逐个几何纤维而言，不是说整个态射总共只有一个分歧点。“二次的”按上面的局部纤维 $\operatorname{Spec}(k'[t]/(t^2))$ 理解；不额外断言该分歧在任意特征下都是温和分歧。底本未在此另给证明，本译文也未独立验证线性束的存在性及推论的全部假设。}

% NEW_TRANSLATION_D17: EGAVaiello.pdf, PDF/printed pp.41--43, section 11 in full.
\egasection{§11\quad Grassmann 概形}
\source{\texttt{EGAVaiello.pdf} 第41--43页；第十七版补译}

现在我们要使用 $P$ 中的线性子簇，其相对维数不再限于0和 $n-1$，因此显然需要引入关于 Grassmann 概形的一些记号，以及关于线性子簇构造的若干属于“初等几何”的连锁事实〔原词：\textit{sorites}〕\ednote{这个词最好怎样翻译？〔底本脚注57。〕}；它们全都应当放在本段开头。此外，实际中有时也要取任意线性截面，而不只取超平面截面，所以宜以这种较广的观点重新审视前面所有编号小节。

设 $E$ 是预概形 $S$ 上的拟凝聚模，$n>0$ 为整数。考虑函子
\[
(\mathbf{Sch}/S)^{\mathrm{op}}\longrightarrow\mathbf{Ens},
\]
其在 $S'$ 上的值 $\operatorname{Grass}_n(E)(S')$，为 $E_{S'}$ 的局部自由、秩为 $n$ 的商模所成的集合。\transnote{这里的商模按相容的商映射同构类计；下标 $S'$ 表示沿 $S'\to S$ 的拉回。底本写 $(\mathrm{Sch})^0/S$ 及 $E'_s$，正文将反变范畴和拉回记号写清。}

这个函子可表；代表它的 $S$-预概形也记作 $\operatorname{Grass}_n(E)$。为证明可表性，考虑函子间的自然态射
\[
\operatorname{Grass}_n(E)\longrightarrow
\operatorname{Grass}_1(\textstyle\bigwedge^nE)
=\mathbf P(\textstyle\bigwedge^nE),
\]
它把 $E_{S'}$ 的每个局部自由、秩为 $n$ 的商模 $G$，对应到局部自由、秩为一的模 $\bigwedge^nG$，将后者看作 $(\bigwedge^nE)_{S'}$ 的商模。\cmt{D17-E04}{底本在说明这条自然态射时，把秩一的模和它的源都仍写为 $G$、$E_{S'}$，遗漏了外幂。秩为 $n$ 的 $G$ 一般不是秩一模；正文补全明确给出此映射的 $\bigwedge^nE_{S'}\twoheadrightarrow\bigwedge^nG$。}
与 Cartan 讨论班\ednote{请把参考文献写得更精确。〔底本脚注58。〕}中的证明一样，可以证明这个态射“可由闭浸入表示”；于是 $\operatorname{Grass}_n(E)$ 是 $\mathbf P(\bigwedge^nE)$ 的闭子概形。特别地，它在 $S$ 上分离；当 $E$ 为有限型时，它在 $S$ 上拟紧并且射影。\cmt{D17-E05}{底本把“在 $S$ 上拟紧”无条件置于“若 $E$ 为有限型则射影”之前。一般拟凝聚 $E$ 不保证拟紧性：取 $S=\Spec k$、$E=\bigoplus_{i\geq0}ke_i$、$n=1$，则 $\mathbf P(E)=\operatorname{Proj}k[T_0,T_1,\ldots]$ 的标准开覆盖 $D_+(T_i)$ 没有有限子覆盖。正文将拟紧性限制在此处足够的有限型情形，详细反例见核查报告。}
若 $E$ 为有限表示，则 $\operatorname{Grass}_n(E)$ 在 $S$ 上也为有限表示。事实上，可以假设 $S$ 仿射，写作 $S=\Spec A$；于是 $E$ 可以由 $A$ 的某个有限型子环上的有限表示模通过扩标量得到，而 $\operatorname{Grass}_n(E)$ 的构造显然与 $S$ 上的基变换相容。\transnote{底本这里只说“来自 $A$ 的某个有限型子环上的有限型模”。为使所述归约真正保持有限表示，取定义 $E$ 的有限表示矩阵的有限个系数所生成的子环，并在该子环上使用同一有限矩阵即可；正文将这项必要信息写明。这里不是断言任意有限型模都能如此下降其全部关系。}

若 $E$ 局部自由且秩有限，则 $\operatorname{Grass}_n(E)$ 在 $S$ 上光滑；在秩至少为 $n$ 的部分，其纤维几何连通。这来自一个更精确的事实：若 $E$ 自由且秩为 $r\geq n$，则 $\operatorname{Grass}_n(E)$ 可以由 $\binom rn$ 个开子集覆盖，每个开子集都在 $S$ 上同构于相对维数为 $n(r-n)$ 的仿射空间。\transnote{底本未说明 $n>r$ 时的边界；此时没有秩为 $n$ 的商，代表概形为空。正文将非空的几何连通纤维断言限制在 $n\leq r$，而空部分仍在 $S$ 上光滑。}
这一覆盖来自于借助 $E$ 的一组基选取 $\binom rn$ 个坐标直和分解，即分裂正合列\transnote{底本在这条正合列首项把 $E'$ 印成 $E$，而次页明说 $E'$ 秩为 $n$；按坐标分解订正。}
\[
(s)\qquad0\longrightarrow E'\longrightarrow E
\longrightarrow E''\longrightarrow0,
\qquad\operatorname{rank}E'=n.
\]

这样的正合列确定 $\operatorname{Grass}_n(E)$ 的一个子函子 $\operatorname{Grass}_n(s)$：只取那些使复合同态
\[
E'_{S'}\longrightarrow E_{S'}\longrightarrow G
\]
为满射、从而为同构的局部自由秩 $n$ 商 $G$。包含态射
\[
\operatorname{Grass}_n(s)\longrightarrow\operatorname{Grass}_n(E)
\]
可由开浸入表示；另一方面，在上述选定的坐标分裂下，$\operatorname{Grass}_n(s)$ 典范同构于向量丛\cmt{D17-E06}{底本在此重复一段“可由开浸入表示；另一方面”，并把表示所选分解的 $s$ 多次写成 $S$。正文删去重复、统一记号。这个开图参数化从 $E$ 到 $E'$ 的、在 $E'$ 上为恒等的同态；在选定分裂后，它由 $E''\to E'$ 给出。在本书 $\mathbf V(M)=\operatorname{Spec}\operatorname{Sym}M$ 的约定下，上式正好代表这种同态。若仅给任意未选分裂的短正合列，则该开图是相应向量丛的挠子，并无无条件的典范零截面；正文保留了坐标分解所给的分裂前提。}
\[
\mathbf V\bigl(\mathscr{H}\!om_{\mathcal O_S}(E',E'')\bigr).
\]


这种特殊结构的一个推论是：若 $s\in S$，则在 $E$ 局部自由、秩有限的条件下，$\operatorname{Grass}_n(E)$ 的每个 $k(s)$-值点都可提升为 $s$ 的某个邻域上的截面。另一方面，若 $S=\Spec k$，$k$ 是无限域，则 $\operatorname{Grass}_n(E)$ 的每个非空开子集都包含一个 $k$-有理点。

$\operatorname{Grass}_n(E)$ 的一个 $S$-值点，即 $E$ 的一个局部自由秩 $n$ 商模 $G$，典范地确定 $\mathbf P(E)$ 的一个子概形 $\mathbf P(G)$。这样的子概形称为 $\mathbf P(E)$ 的\emph{线性子簇}（可能混淆时，说相对于 $S$ 的线性子簇）；这里在定义名称时不限定或指定 $G$ 的秩〔英译者说明〕。若 $n\geq1$，它是相对维数 $n-1$ 的射影纤维丛；若 $n=0$，它为空。立即可以验证，若知道对应的线性子簇，就知道 $\operatorname{Grass}_n(E)$ 的截面，也就是商 $G$。这样，对于 $n\geq1$，Grassmann 概形可解释为函子“$\mathbf P(E)_{S'}$ 的相对维数 $n-1$ 的线性子簇”的代表，其中 $S'$ 变动。

此外，还能给出后一函子、也就是相对维数 $m$ 的线性子簇概念的内在刻画：它们在每个 $s\in S$ 处“射影次数为一”。这一刻画将在后面的章节中给出，我们这里完全用不到它。\transnote{底本这句语法有缺口，只写出相对维数与逐纤维“射影次数一”，没有完整列明平坦性、有限表示等可能需要的条件。这里保留作者的预告，不将该残句单独表述为一个已经完备的充要判据。}

再假设 $E$ 局部自由、秩为 $r$，并以 $E^\vee$ 表示其对偶。由极性对应，得到典范同构
\[
\operatorname{Grass}_n(E)\simeq
\operatorname{Grass}_{r-n}(E^\vee),
\]
它把 $E$ 的商 $G$ 对应到 $E^\vee$ 的商 $E^\vee/G^\vee$。从线性子簇的观点看，$P$ 中相对维数为 $m$ 的线性子簇 $L$，对应到 $P^\vee$ 中相对维数为
\[
r-m-2=(r-1)-1-m
\]
的对偶线性子簇 $L^0$，也就是说，其在 $P^\vee$ 中的相对余维为 $m+1$。〔附记：这里 $r-1$ 是 $P$ 和 $P^\vee$ 在 $S$ 上的共同相对维数。〕\transnote{底本将维数印作 $(r-1)+-1-m$，正文只整理这两个相邻运算符。$m=r-1$ 时对偶线性子簇为空，依本段的约定赋予维数 $-1$。}
几何上可以如下理解这个对应。先取 $n=r-1$，便得到同构
\[
\mathbf P(E^\vee)\simeq\operatorname{Grass}_{r-1}(E),
\]
从而可将 $P^\vee$ 的 $S$-值点（例如取 $S$-值点）认同为 $P$ 中余维为一的线性子簇，仍称为 $P$ 的\emph{超平面}。

这说明 $L^0$ 由包含 $P$ 中线性子簇 $L$ 的超平面组成；精确地说，$L^0$ 的 $S'$-值点就是 $P_{S'}$ 中包含 $L_{S'}$ 的超平面。\transnote{底本在括号中写成“$L$ 的 $S'$-值点”，但所描述的是对偶参数空间 $L^0$ 的点；正文补上上标。}
这来自下面的事实；它本应与“$P$ 的线性子簇 $L$ 确定其所来自的局部自由商 $G$（或 $Q$，字迹不清）”这一事实同时出现：若 $G$ 与 $G'$ 是 $E$ 的两个局部自由商（不一定同秩），则
\[
\mathbf P(G')\subset\mathbf P(G)
\]
作为 $\mathbf P(E)$ 的线性子簇之间的包含，等价于 $G'$ 是 $G$ 的商；这个包含正是由 $G\twoheadrightarrow G'$ 诱导的态射。

以上是我们至少应当掌握的连锁事实〔法文：\textit{sorites}〕\ednote{请询问 A.G. 或 Deligne，最好用哪个词。〔底本脚注59。〕}。在本段\ednote{A.G. 所说的 paragraph 是什么意思？〔底本脚注60。〕}所有其他编号小节写出之前，无论如何还不能确定完整的清单。

我觉得也宜引入函子 $\operatorname{Grass}(E)$，其在 $S'$ 上的值为 $E_{S'}$ 的局部自由、有限秩但不预先指定秩的商模所成的集合〔字迹不清，请询问 A.G.〕。于是它由
\[
\coprod_{n\geq0}\operatorname{Grass}_n(E)
\]
代表。$\mathbf P(E)$ 的线性子簇确实由 $\operatorname{Grass}(E)$ 在 $S$ 上的截面确定。〔附记：若 $S$ 不连通，秩、也就是相对维数，可以变化。此处字迹略有不清，请与 A.G. 核实。〕\transnote{这里“秩不预先指定”指局部有限秩，秩函数可在不同开闭部分取不同的非负整数；不是允许无限秩商。否则上面的有限整数指标不交并不能代表那个更大的函子。}

